% !TEX root = ../main.tex
\chapter{隐函数定理与比较静态}
\label{ch:ift}

\noindent\emph{用到的前置工具：偏导数、梯度、Jacobian 与 Hessian（第十三章）；矩阵求逆与正定矩阵（线性代数部分，正定性见第五章）。}
\medskip

如果要在经济学家天天使用的数学工具里只选一件作为``镇店之宝'',隐函数定理很可能当选。它之所以重要，不是因为它有多深，而是因为经济学问的问题恰好是它最擅长回答的那一类。本章先点明这类问题长什么样，再把单变量与多变量两个版本讲清楚，最后用两个经典场景说明：从市场的税负转嫁，到由一阶条件做的比较静态，背后都是同一把工具在起作用——甚至连``二阶条件''这种看似只关乎``是不是极大值''的东西，真正让比较静态有良定义的也是它。

\section{从``解不出来''说起}
\label{sec:ift-motivation}

经济学的结论，绝大多数是\textbf{比较静态}（comparative statics）的结论：某个外生条件——一项税、一笔收入、一个成本参数——动一下，作为均衡或最优选择的内生量会朝哪个方向、变动多少？需求曲线右移，价格涨还是跌、涨多少？征一笔从量税，消费者和生产者各承担多大份额？这些问句的数学形式，统统是``某个内生量对某个参数的导数''。

把这类问题抽象出来，总是同一个结构：有一个把内生量 $x$ 与参数 $p$ 联系起来的方程
\[
    F(x,p)=0 ,
\]
它可能是``供给等于需求''的市场出清条件，可能是``边际收益等于边际成本''的最优化一阶条件，也可能是某个不动点方程。我们真正关心的，是由它\emph{隐含地}定义出来的关系 $x=x^*(p)$，尤其是它的导数 $\dd{x^*}{p}$——这正是比较静态。

麻烦在于：上面这个方程通常\textbf{解不出显式的 $x^*(p)$}。需求与供给可能是任意的非线性函数，一阶条件可能是超越方程。换句话说，我们想要导数 $\dd{x^*}{p}$，却连函数 $x^*(p)$ 本身都写不出来。

隐函数定理正是为此而生。它告诉我们：只要在某一点满足一个温和的``非退化''条件，就\emph{不必解出} $x^*(p)$，也能断定这个函数在该点附近确实存在、而且可微，并直接给出它的导数表达式。这正是``工具性''的极致——把一个看似要先求解才能回答的问题，变成一道纯粹的求导练习。

\section{单变量情形}
\label{sec:ift-univariate}

\thm{隐函数定理（单变量）}{
设 $F:\RR^2\to\RR$ 在点 $(x_0,p_0)$ 的某个邻域内连续可微，且
\[
    F(x_0,p_0)=0,
    \qquad
    \pd{F}{x}(x_0,p_0)\neq 0 .
\]
则存在 $p_0$ 的一个邻域 $U$ 与定义在 $U$ 上的连续可微函数 $x=x^*(p)$，使得 $x^*(p_0)=x_0$，并且对一切 $p\in U$ 都有 $F\pr{x^*(p),p}=0$。此外，
\[
    \boxed{\ \dd{x^*}{p}
    = -\,\frac{\partial F/\partial p}{\partial F/\partial x}\ } .
\]
}

定理分两部分：``存在且可微''是较深的部分，``导数等于这个比值''是我们日常真正使用的部分。后者其实一眼可得。既然在 $U$ 上恒有
\[
    F\bigl(x^*(p),\,p\bigr)\equiv 0,
\]
这是一个关于 $p$ 的\emph{恒等式}，两边对 $p$ 求导（左边用链式法则，把 $x^*$ 看成 $p$ 的函数）便得
\[
    \pd{F}{x}\cdot\dd{x^*}{p}+\pd{F}{p}=0 .
\]
只要 $\partial F/\partial x\neq 0$，就能把 $\dd{x^*}{p}$ 解出来，正是定理中的公式。可以看到，那个看似技术性的条件 $\partial F/\partial x\neq 0$，恰恰就是``这一步除法做得下去''的保证。

\rmkb[存在性是更深的部分]{
``$x^*(p)$ 存在且可微''这半句话，才是隐函数定理真正的数学内容，其严格证明要用到压缩映射或反函数定理（见第十二章）。直觉上它说的是：曲线 $F=0$ 在该点附近``不打弯成竖直'',因而可以被看成以 $p$ 为自变量的一条函数图像。本书按使用者的需要，把存在性当作已知、把重点放在导数公式与它的经济学含义上。
}

\paragraph{几何图像。}
方程 $F(x,p)=0$ 在 $(p,x)$ 平面上描出一条曲线——它是 $F$ 的一条等值线（零水平集）。条件 $\partial F/\partial x\neq 0$ 意味着这条曲线在该点处\emph{不是竖直的}，于是在该点附近，曲线可以写成 $x$ 关于 $p$ 的函数图像 $x=x^*(p)$（图~\ref{fig:ift-levelset}）。沿着这条曲线移动时，必须保持 $F$ 的值不变（始终为 $0$），全微分 $\partial F/\partial x\,\d x+\partial F/\partial p\,\d p=0$ 表达的正是这件事，而它解出的 $\d x/\d p$ 就是曲线的斜率。

\begin{figure}[h]
\centering
\begin{tikzpicture}[scale=1.05,>=Stealth]
    \draw[->] (-0.2,0) -- (6.4,0) node[right] {$p$};
    \draw[->] (0,-0.2) -- (0,4.2) node[above] {$x$};
    % 曲线 F(x,p)=0，即局部的函数图像 x=x^*(p)
    \draw[thick,domain=0.35:5.8,smooth,variable=\t] plot ({\t},{1.2+0.9*sqrt(\t)});
    % 曲线标签：置于曲线右端的上方空白处，不压线
    \node[anchor=west] at (4.55,3.55) {$F(x,p)=0$};
    % 切点（精确落在曲线上：x=1.2+0.9*sqrt(3)=2.76）
    \coordinate (P) at (3,2.76);
    \fill (P) circle (1.6pt);
    % 辅助虚线引到两轴，标签置于轴外
    \draw[dashed] (P) -- (3,0) node[below] {$p_0$};
    \draw[dashed] (P) -- (0,2.76) node[left] {$x_0$};
    % 切线段（斜率 0.26），贴合曲线于 P
    \draw[thick,blue!60!black] (2.1,2.526) -- (3.9,2.994);
    % 斜率标签停在左上大片空白，细引线指向切线（全程在曲线上方，不交叉）
    \node[blue!60!black,anchor=west] at (0.5,3.6) {切线斜率 $=\dfrac{\d x^*}{\d p}$};
    \draw[blue!60!black,->] (2.05,3.35) -- (2.55,2.70);
\end{tikzpicture}
\caption{零水平集 $F(x,p)=0$ 在 $(x_0,p_0)$ 附近不竖直，故可表为函数图像 $x=x^*(p)$；其斜率即由隐函数定理给出。}
\label{fig:ift-levelset}
\end{figure}

\rmkb[为什么``$\partial F/\partial x\neq 0$''不可或缺]{
看 $F(x,p)=x^2-p$ 在原点 $(0,0)$：此处 $\partial F/\partial x=2x=0$，条件失效。事实上 $F=0$ 给出 $x=\pm\sqrt{p}$——对 $p>0$ 有上下两支、对 $p<0$ 无解，原点附近 $x$ 根本不是 $p$ 的单值函数（曲线在此处切线竖直）。这时公式里的分母为零、导数``爆掉'',对应的经济含义往往是均衡分叉或多重均衡。换言之，那个非退化条件不是技术性累赘，它在守护``存在唯一一条可微的内生关系''这件事本身。
}

\section{应用一：税负转嫁}
\label{sec:ift-tax}

考虑一个单一市场：消费者按价格 $p$ 需求 $D(p)$，$D'<0$；政府对每单位交易向生产者征收从量税 $t$，于是生产者实得 $p-t$，供给为 $S(p-t)$，$S'>0$。市场出清要求需求等于供给。把均衡价格 $p$ 当内生量、税率 $t$ 当参数，出清条件写成
\[
    F(p,t)\;=\;D(p)-S(p-t)\;=\;0 .
\]
我们想知道 $\dd{p^*}{t}$：征税后消费者支付的价格会涨多少？这正是``税最终由谁承担''的问题。即便 $D,S$ 形式未知、$p^*(t)$ 解不出来，隐函数定理依然给得出答案。先算两个偏导：
\[
    \pd{F}{p}=D'(p)-S'(p-t)<0,
    \qquad
    \pd{F}{t}=S'(p-t)>0 .
\]
于是
\[
    \dd{p^*}{t}
    =-\,\frac{\partial F/\partial t}{\partial F/\partial p}
    =-\,\frac{S'}{D'-S'}
    =\frac{S'}{\,S'-D'\,}\;\in(0,1).
\]

这个简洁的式子蕴含了全部的经济学。分母 $S'-D'>0$，故 $\dd{p^*}{t}\in(0,1)$：\textbf{消费者价格的涨幅小于税额本身}，涨的那部分 $\dfrac{S'}{S'-D'}$ 就是\emph{消费者承担的税负份额}，余下的 $\dfrac{-D'}{S'-D'}$ 由生产者承担。当供给极富弹性（$S'\to\infty$）时份额趋于 $1$，税几乎全压给消费者；当需求极富弹性（$D'\to-\infty$）时份额趋于 $0$，生产者几乎全担——``谁更难调整，谁就承担更多''。这一整套结论，我们没有解出 $p^*(t)$ 就拿到了。

\section{多变量情形}
\label{sec:ift-multivariate}

经济学里内生量往往不止一个：一般均衡有一组价格，最优化有一组选择变量。把单变量定理升级到向量，关键的``非退化''条件从``一个偏导非零''变成``一个矩阵可逆''。

\thm{隐函数定理（一般形式）}{
设 $\mat F:\RR^{n}\times\RR^{m}\to\RR^{n}$，记作 $\mat F(\vc x,\vc p)$，其中 $\vc x\in\RR^{n}$ 为内生变量、$\vc p\in\RR^{m}$ 为参数。设 $\mat F$ 在 $(\vc x_0,\vc p_0)$ 邻域内连续可微，$\mat F(\vc x_0,\vc p_0)=\vc 0$，且关于内生变量的 Jacobian
\[
    \D_{\vc x}\mat F=\br{\pd{F_i}{x_j}}_{n\times n}
\]
在该点\emph{可逆}。则在 $\vc p_0$ 附近存在唯一的连续可微映射 $\vc x=\vc x^*(\vc p)$，满足 $\vc x^*(\vc p_0)=\vc x_0$ 与 $\mat F\bigl(\vc x^*(\vc p),\vc p\bigr)\equiv\vc 0$，并且
\[
    \boxed{\ \D_{\vc p}\vc x^*
    =-\,\br{\D_{\vc x}\mat F}^{-1}\,\D_{\vc p}\mat F\ } .
\]
}

推导与单变量如出一辙。把恒等式 $\mat F\bigl(\vc x^*(\vc p),\vc p\bigr)\equiv\vc 0$ 两边关于 $\vc p$ 做全微分（链式法则的矩阵形式）：
\[
    \D_{\vc x}\mat F\;\D_{\vc p}\vc x^*+\D_{\vc p}\mat F=\mat 0 .
\]
因为 $\D_{\vc x}\mat F$ 可逆，左乘其逆即解出 $\D_{\vc p}\vc x^*$。单变量里那个标量 $\partial F/\partial x\neq 0$，恰好是 $n=1$ 时``矩阵可逆''的特例——可逆性扮演的就是``这一步矩阵除法做得下去''的角色。

\section{应用二：从一阶条件做比较静态}
\label{sec:ift-foc}

这是隐函数定理在经济学里最常见、也最值得记住的用法。设决策者求解
\[
    \max_{\vc x\in\RR^{n}}\; f(\vc x;\vc\theta),
\]
$\vc x$ 是选择变量、$\vc\theta\in\RR^m$ 是参数（偏好、价格、禀赋……）。内点最优解满足一阶条件
\[
    \grad_{\vc x} f(\vc x^*,\vc\theta)=\vc 0 .
\]
这恰好是一个形如 $\mat F(\vc x,\vc\theta)=\vc 0$ 的方程组，其中 $\mat F=\grad_{\vc x} f$。对它套用隐函数定理，关键的 Jacobian 是
\[
    \D_{\vc x}\mat F=\D_{\vc x}\grad_{\vc x} f=\grad^2_{\vc x\vc x} f=\mat H,
\]
也就是 $f$ 关于选择变量的 \textbf{Hessian}——正是二阶条件里那个矩阵。于是比较静态公式为
\[
    \D_{\vc\theta}\vc x^*=-\,\mat H^{-1}\,\grad^2_{\vc x\vc\theta} f,
    \qquad
    \grad^2_{\vc x\vc\theta} f=\br{\pdc{f}{x_i}{\theta_k}} .
\]

\state{二阶条件不只是``判极值'',它还让比较静态有良定义}{
内点极大值的二阶充分条件是 Hessian $\mat H$ \emph{负定}。负定矩阵必然可逆——而这恰恰就是隐函数定理对 $\mat F=\grad_{\vc x}f$ 的适用前提！于是：只要你找到的是一个满足二阶条件的内点极大值，比较静态 $\D_{\vc\theta}\vc x^*=-\mat H^{-1}\grad^2_{\vc x\vc\theta} f$ 就\emph{自动}有定义、并可由此算出。换句话说，二阶条件一身二任：既保证你停在山顶，又保证``参数动一下、最优选择平滑地跟着动''。
}

在标量情形（一个选择变量 $x$、一个参数 $\theta$）这条公式格外透亮：
\[
    \dd{x^*}{\theta}=-\,\frac{\partial^2 f/\partial x\partial\theta}{\partial^2 f/\partial x^2}
    =-\,\frac{f_{x\theta}}{f_{xx}} .
\]
二阶条件是 $f_{xx}<0$，所以
\[
    \sgn\!\pr{\dd{x^*}{\theta}}=\sgn\!\pr{f_{x\theta}} .
\]
\textbf{最优选择随参数变动的方向，完全由目标函数的交叉偏导符号决定}：若 $\theta$ 提高了选择变量的边际收益（$f_{x\theta}>0$，二者``互补''），则最优 $x^*$ 随 $\theta$ 上升，反之下降。这个朴素的洞见，正是``单调比较静态''（第二十三章，Topkis 与超模性）的种子——后者更进一步，连可微、连 Hessian 可逆都不再需要，只靠``互补''这个序结构就能定出方向。

\ex[需求截距上升如何抬高垄断产量]{
线性需求 $p=a-bq$（$a,b>0$）、单位成本 $c$ 的垄断者，利润为 $\pi(q;a)=(a-bq)q-cq$。求需求截距 $a$ 上升对最优产量的影响 $\dd{q^*}{a}$。
}
\sol{
这里 $f=\pi$、选择变量是 $q$、参数取 $a$。所需的两个二阶偏导为
\[
    \pi_{qq}=\pdn{2}{\pi}{q}=-2b<0\quad(\text{二阶条件满足}),
    \qquad
    \pi_{qa}=\pdc{\pi}{q}{a}=\pd{}{a}\pr{a-2bq-c}=1 .
\]
代入公式：
\[
    \dd{q^*}{a}=-\frac{\pi_{qa}}{\pi_{qq}}=-\frac{1}{-2b}=\frac{1}{2b}>0 .
\]
需求截距上升，最优产量随之上升。本例的一阶条件 $a-2bq-c=0$ 能直接解出 $q^*=(a-c)/(2b)$，对 $a$ 求导同样得 $1/(2b)$——正好用这个能解的简单情形，印证隐函数定理给出的公式无误。真正的价值在于：当一阶条件无法显式求解时，上面这套做法照旧成立。
}

\section{失效的时候，以及它通向何处}
\label{sec:ift-beyond}

\state{隐函数定理失效意味着什么}{
当 $\D_{\vc x}\mat F$ 奇异（行列式为零、Hessian 退化）时定理不适用。在经济学里这通常对应：均衡分叉、多重均衡、或目标函数恰在某个参数值上从单峰变为多峰的临界点。此时比较静态可能不连续、甚至无定义。两条出路：(i) 退一步，只问``方向''而不问``大小''——这正是\textbf{单调比较静态}（第二十三章）的领域，它不要求可微、也不要求 Hessian 可逆，只依赖目标函数的互补（超模）结构；(ii) 借助更强的拓扑工具（不动点、度理论，第十一章）讨论均衡的存在与个数。
}

最后值得记住的是，本章这把工具远不止用于微观。它的几处重要去处：

\begin{itemize}
    \item \textbf{包络定理}（第二十二章）补上了另一半故事——隐函数定理告诉你``最优\emph{选择}如何随参数变''，包络定理告诉你``最优\emph{值}如何随参数变'',二者常常配套出现。
    \item \textbf{Slutsky 方程}是消费者理论里隐函数定理的集大成应用：对马歇尔需求背后的一阶条件做隐函数微分，就分解出替代效应与收入效应。
    \item \textbf{估计量的渐近分布}（第三十四章）。M-估计量 $\hat{\vc\theta}$ 由样本一阶条件 $\frac{1}{n}\sum_i \grad_{\vc\theta}\,m(\vc z_i,\hat{\vc\theta})=\vc 0$ 隐含定义；把它对样本扰动做隐函数式的展开，正是导出 $\rtn(\hat{\vc\theta}-\vc\theta_0)\dto\cN(\cdot)$ 的核心一步——计量里的 Delta 法与三明治协方差矩阵，骨子里都是隐函数定理。
\end{itemize}

一件工具，三门核心课。这正是本书想反复呈现的：少数几个数学构造，被反复用在看似无关的经济学问题上；认清了工具本身，就等于同时拿到了许多扇门的钥匙。
